\section{Exploratory Data Analysis}

\subsection{Dataset overview and data model}

\code{broadway\_clean.csv} contains 29{,}167 rows and 12 columns at a grain of
one row per \emph{week $\times$ show $\times$ theatre}.\footnote{%
\textbf{On tables and joins.} The brief asks for the tables used and joins
performed. This dataset is a \emph{single denormalised file}: no joins are
required and none were performed. We meet the underlying obligation ---
documenting the data model --- by recovering the entity structure the flat file
implies (Figure~\ref{fig:datamodel}).}
The triple (\code{Date.Full}, \code{Show.Name}, \code{Show.Theatre}) is a
verified natural key with no duplicates, and every observation is a week ending
on a Sunday, the Broadway League's reporting convention. The file spans 1990-08-26 to 2016-08-14
with 1{,}281 distinct week-endings, 820 shows and 56 theatres, and a median of
25 productions reporting per week.

Although denormalised, the data encodes a three-level hierarchy
(Figure~\ref{fig:datamodel}): \textbf{theatres} (fixed seat count),
\textbf{shows} (the creative work), and \textbf{productions} --- one show's run
in one theatre, the unit at which commercial decisions are taken. That
distinction matters, since 27 titles appear in more than one theatre and some
return after a break, so we key a production on
(\code{show}, \code{theatre}, \code{run\_block}), incrementing \code{run\_block}
after any absence beyond four weeks. This yields 942 productions overall, 902
inside the analysis window; Table~\ref{tab:dictionary} is the full data
dictionary.

\subsection{Data quality audit}

The file has no nulls, no duplicate rows and no formatting faults. That apparent
cleanliness is misleading: missingness has been encoded as zero, and two further
defects are structural. Table~\ref{tab:quality} records every issue, its
evidence and our decision; the decisions live in \code{src/data\_load.py}, so
code and table cannot drift apart. In summary: 2{,}309 rows (7.9\%) report zero
performances alongside positive attendance, which is impossible; 1{,}918 rows
(6.6\%) code gross potential as zero, including 96\% of all 1996 rows; capacity
is censored at 100\% while 1{,}433 rows show gross potential above it under
premium pricing; and 75 weeks are missing outright, almost all before 1993.

Critically, the zero-coded values cluster by \emph{year} rather than by show
(Figure~\ref{fig:quality}b), identifying them as reporting artefacts of
particular eras. Together with the thin early coverage --- a handful of
productions per week before 1995 against roughly 25 afterwards --- this is what
drives our decision to begin the analysis in 1996, leaving 28{,}479 rows.

\paragraph{Recovering an undocumented column.}
\code{Statistics.Capacity} is named as a seat count but ranges only from 10 to
100. If it is instead the percentage of seats sold, then
\begin{equation}
\text{seats} \;=\; \frac{\text{attendance}}{\text{performances}}
\Big/ \frac{\text{capacity}}{100}
\label{eq:identity}
\end{equation}
must recover a constant per theatre. It does: implied values are tight within
each venue and match published seat counts to a median error of 1.7\%, with the
Lunt-Fontanne recovering exactly its published 1{,}509 seats
(Figure~\ref{fig:validation}). This settles the column's meaning and supplies a
legitimate static feature in place of a 56-level categorical.

\subsection{Temporal patterns}

Total weekly gross grows roughly four-fold across the series
(Figure~\ref{fig:market}), but almost none of that is demand. Between 1996 and
2015 the median ticket price rises from \$44.60 to \$93.10 (+109\%) while median
capacity utilisation moves only from 82\% to 85\% (Figure~\ref{fig:price}). A
model trained on dollar levels would spend most of its capacity learning the
price trend --- hence our scale-free targets.

Demand is strongly seasonal, but \emph{only at week resolution}. ISO weeks
52--53 reach a median of 90\% capacity against 73\% in weeks 36 and 44, a
17-point swing (Figure~\ref{fig:seasonality}). Aggregated to months that
contrast is exactly zero: December's median equals September's, because the
holiday peak is a two-week phenomenon cancelled by the slow first three weeks of
December. Calendar features must therefore be built at week resolution, and ours
are.

Two weeks are dominated by exogenous shocks: in the week ending 16 September
2001 market gross fell from \$7.48m to \$3.47m ($-54\%$), and during the
November 2007 stagehands' strike\footnote{%
\textbf{Note.} A labour dispute between the stagehands' union and the League of
American Theatres and Producers closed most Broadway houses for nineteen days
in November 2007. Only the handful of productions outside the affected
contracts continued to play, which is why the count of reporting productions
--- not merely their revenue --- collapses in Figure~\ref{fig:market}b.}
the productions able to play fell from 33 to 8 while gross collapsed from
\$17.1m to \$4.6m. Neither is predictable from
anything in this dataset. We retain and flag both rather than deleting
inconvenient rows, and treat them as evidence that a real share of the variance
here is irreducible.

\subsection{Shows, theatres and lifecycles}

Of the brief's three distribution analyses, two map onto this data: the
\emph{product} is the show, the \emph{seller} is the theatre. The third cannot
be done --- the file records aggregate weekly attendance, never an individual
buyer --- so no customer-level problem is scopeable in Phase 2.

Show types are heavily imbalanced --- 71.1\% musicals, 27.8\% plays, 1.1\%
\code{Special} (Figure~\ref{fig:types}) --- so a classifier that always answers
``Musical'' scores 71\% while being useless, illustrating the metric problem of
Section~1.5. Recovered seat counts span roughly 500 to 1{,}900 and are close to
independent of utilisation (Figure~\ref{fig:theatres}), so both carry
information.

Run length is the most skewed quantity in the data: median 15 weeks, 90th
percentile 68, maximum 994 (\emph{The Phantom of the Opera} at the Majestic).
Twenty-four productions were still running when the data ends and are therefore
right-censored. Kaplan--Meier estimates \citep{kaplan1958nonparametric} handling
that censoring give median survival of 20 weeks for musicals against 13 for
plays (Figure~\ref{fig:survival}), reproducing the musical-longevity effect of
\citet{simonoff2016survival}.

\subsection{Relationships and the leakage they imply}

Capacity is not an independent measurement but an identity over attendance,
performances and seat count (Equation~\ref{eq:identity}). Inverting it
reconstructs capacity from same-week columns to a median absolute error of
\textbf{0.48 percentage points}, so a model given same-week attendance or gross
is not forecasting demand --- it is reading the answer. This is the dominant
leakage risk here \citep{kaufman2012leakage}, and we enforce it in code: a guard
rejects any feature list containing a same-week outcome column.

\subsection{Candidate problem exploration}
\label{sec:candidates}

We examined three candidates. All probes use plain logistic or linear regression
inside a \code{scikit-learn} pipeline, deliberately untuned, on a chronological
split (train $\le$ 2013, validation 2014, test 2015 onward) that also scores the
90 test productions never seen in training.\footnote{%
\textbf{Note on splitting.} The brief recommends stratified splits for
imbalanced classification. Chronology takes precedence here: a stratified
random split would train on 2016 and test on 2004, and would break the
production-grouping constraint. Our three-band target is balanced by
construction, so stratification buys little; for the imbalanced Candidate 2
target we use class weights and time-ordered folds instead.}
A random split would be badly optimistic: \emph{Phantom} alone contributes 994
rows, so splitting by row grades the model on interpolating a run it has
memorised \citep{bergmeir2012cv}. This is the Broadway analogue of the brief's
\code{customer\_id} warning --- the unit that must not straddle the split is the
production.

\paragraph{Candidate 1 --- demand for a running production (primary).}
One question framed both ways: regression on capacity utilisation, and
classification into three bands at training-set terciles (\emph{Low} $<73\%$,
\emph{Mid}, \emph{High} $\ge 89\%$). The stakeholder is a theatre general
manager setting staffing, discounting and TKTS allocation.\footnote{%
\textbf{Note.} TKTS is the Times Square discount booth run by the Theatre
Development Fund, where producers release unsold seats at a reduced price. How
many to release, and how far ahead, is precisely the decision a demand forecast
supports --- and it is committed weeks in advance, not on the day.}

The horizon is a scoping decision, made with evidence
(Table~\ref{tab:horizon}). Capacity is highly persistent --- lag-1
autocorrelation 0.84 --- so one week ahead a naive ``same as last week'' rule
scores 0.728 and the model cannot beat it; but a one-week forecast is also
useless, since marketing commitments cannot be changed at that notice. At four
weeks persistence has decayed to 0.607 while the model holds 0.638, and the
regression cuts MAE by 8.1\%. We therefore scope Candidate 1 at a
\textbf{four-week horizon}, where the forecast is both actionable and
non-trivial.

There the classifier reaches \textbf{0.638} accuracy against a 0.315 majority
and 0.607 persistence baseline, with balanced accuracy 0.637 and 0.651 on unseen
productions; the regression reaches MAE 7.36 and $R^2$ 0.546 against 8.01 and
0.399. Errors concentrate in the \emph{Mid} band --- the expected artefact of
cutting a continuous quantity at tercile boundaries, not a modelling failure.

The guard's value is measurable: adding same-week attendance lifts accuracy from
0.638 to 0.875, and adding performances too to 0.902. Those are not
achievements; they are what this dataset's failure mode looks like.

\paragraph{Candidate 2 --- closure risk within eight weeks (secondary).}
The positive rate is 22.1\%, and 208 rows are right-censored because the data
ends before a closure could be observed; those are labelled missing rather than
defaulted to ``no''. With balanced class weights the model attains PR-AUC 0.698
against a no-skill floor of 0.204, ROC-AUC 0.887 and recall 0.806 on the
minority class. Its accuracy of 0.813 barely beats the 0.796 from always
predicting ``stays open'' --- a rule catching none of the closures it exists to
find, which is why PR-AUC is the honest summary \citep{saito2015prplot}. This
candidate contributes the imbalanced, censored target Candidate 1 lacks.

\paragraph{Candidate 3 --- opening-month gross (deprioritised).}
Predicting gross in a production's first four weeks explains only 29\% of the
variance in log gross, essentially all of it theatre size and seasonality.
Without lags little remains, and the variables the literature calls decisive ---
advance sales, casting, marketing spend, reviews \citep{reddy1998broadway} ---
are absent here. We record the boundary rather than force a weak model;
Table~\ref{tab:checklist} scores all three against the brief's checklist.

\subsection{Key insights}

\begin{enumerate}\itemsep1pt
\item A file with zero nulls is not a clean file: 7.9\% of \code{Performances}
and 6.6\% of \code{Gross Potential} are missing data encoded as zero, clustered
by reporting era.
\item \code{Statistics.Capacity} is an undocumented percentage, recoverable from
the data itself to within 1.7\% of published seat counts.
\item Nominal revenue growth is a price effect: ticket prices rise 109\% from
1996 to 2015 while utilisation moves three points.
\item Seasonality exists only at week resolution; monthly aggregation destroys
it entirely.
\item Capacity is reconstructible from same-week columns to 0.48 percentage
points, making leakage the central risk and motivating a guard enforced in code.
\item Forecast difficulty is a design choice: at one week persistence solves the
problem, at four weeks it is genuinely open --- and four weeks is when the
decision is made.
\end{enumerate}

We expect a tuned Candidate 1 classifier to reach the low-to-mid 60s in
accuracy --- where the ceiling honestly sits, given the drivers the literature
identifies are unobserved here. A result above 90\% would indicate leakage, not
skill.
